% NOTE: Filename is fixed by the book outline; content teaches Week 20
% (Unstructured Data & AI Services) of the syllabus, plus the Milestone 5 capstone.
\chapter{Unstructured Data and Managed AI Services}
\index{Amazon Textract}\index{Amazon Rekognition}\index{Amazon Comprehend}\index{Amazon Translate}\index{Amazon OpenSearch}\index{OCR}\index{unstructured data}%
\begin{objectives}
\begin{itemize}
    \item Extract text, key--value pairs, and tables with Amazon Textract, and label image/video content with Amazon Rekognition.
    \item Apply Amazon Comprehend (NLP) and Amazon Translate to document text at pipeline scale.
    \item Integrate AI APIs into Lambda and Glue pipelines with the synchronous/asynchronous boundary respected.
    \item Index extracted metadata into OpenSearch for discovery.
    \item Manage PII risk in persisted OCR output.
    \item Architect a terabyte-scale image pipeline: OCR, object tagging, searchable metadata---and deliver the Milestone~5 end-to-end MLOps capstone.
\end{itemize}
\end{objectives}

\begin{crosscloud}
Pre-trained AI APIs---vision, OCR, NLP, translation---are commodity building blocks on every cloud; the skill is pipeline integration, not model building.

\vspace{2mm}
{\footnotesize
\begin{tabular}{@{}p{2.8cm}p{3.0cm}p{3.0cm}p{3.0cm}@{}}
\toprule
\textbf{Capability} & \textbf{AWS} & \textbf{Azure} & \textbf{GCP} \\
\midrule
OCR / document parsing & Textract & Document Intelligence & Document AI \\
Image/video labeling & Rekognition & Azure AI Vision & Vision AI / Video AI \\
NLP (entities, sentiment, PII) & Comprehend & Azure AI Language & Natural Language AI \\
Translation & Translate & Translator & Translation API \\
Search over metadata & OpenSearch & AI Search & Vertex AI Search \\
Async document jobs & Textract async APIs & Batch APIs & Batch processing \\
\addlinespace
PII detection & Comprehend PII + Macie & Language PII + Purview & Sensitive Data Protection \\
Custom models & Rekognition Custom Labels, Comprehend custom & Custom Vision & Vertex custom training \\
\bottomrule
\end{tabular}}

\vspace{1mm}
\textbf{Fabric} exposes Azure AI services natively in notebooks and pipelines. \textbf{Snowflake} offers Document AI and Cortex functions---extraction as SQL. \textbf{MongoDB} stores the extracted JSON naturally; its vector search serves the retrieval side of the same estate.
\end{crosscloud}

\section{Week 20 Roadmap and Mental Model}
So far, ``data'' meant rows. Week~20 opens the other 80\% of enterprise information: scanned forms, PDFs, images, video, free text. The week's services are \emph{pre-trained APIs}---you do not train Textract or Rekognition; you integrate them. That shifts the engineering content entirely to the pipeline: event-driven invocation, sync/async boundaries, cost per thousand pages, error handling and retries, result persistence, indexing, and the privacy posture of everything the APIs emit \cite{awsTextractDocs,awsRekognitionDocs}.

The mental model: an AI API is a \emph{per-record enrichment transform} with a network boundary. Everything learned in Part~4 about Lambda transforms applies; everything in Part~2 about catalogs and governance applies to the outputs.

\begin{definitionbox}
\textbf{Managed AI services} are pre-trained, API-accessible models---vision, OCR, NLP, translation---billed per unit of input (page, image, character block). They convert unstructured content into structured, catalog-able, indexable data without any model training.
\end{definitionbox}

\begin{keyterms}
\textbf{Textract}: OCR plus structure---key--value pairs, tables, forms; synchronous for single pages, asynchronous for multi-page documents. \textbf{Rekognition}: label, face, text, and moderation detection for images and video. \textbf{Comprehend}: entities, sentiment, key phrases, PII detection, custom classification. \textbf{Translate}: neural machine translation. \textbf{OpenSearch}: managed search/analytics engine for the extracted metadata. \textbf{Confidence score}: every API result carries one; thresholds are a pipeline design input.
\end{keyterms}

\section{Monday: Rekognition and Textract}
\subsection{Vision as an API}
Rekognition answers ``what is in this image/video?'': labels with confidence scores, detected text, faces (analysis, and recognition against a collection you manage), and content moderation categories \cite{awsRekognitionDocs}. For data pipelines the pattern is uniform: S3 object in, JSON result out, confidence thresholds applied downstream.

\subsection{Textract: Structure from Documents}
Textract goes beyond raw OCR: \texttt{AnalyzeDocument} returns forms (key--value pairs) and tables as structured blocks, and \texttt{DetectDocumentText} returns lines and words \cite{awsTextractDocs}. The critical operational boundary: the synchronous APIs accept single-page documents up to a few MB; multi-page PDFs and TIFFs require the \emph{asynchronous} API (\texttt{StartDocumentAnalysis}), which runs a job and publishes completion to SNS. Every document pipeline of consequence is therefore asynchronous---which means it is an event pipeline (S3 upload $\rightarrow$ job $\rightarrow$ SNS $\rightarrow$ result handling), and Chapter~12's orchestration discipline applies.
\index{Textract}\index{asynchronous jobs}

\begin{pitfall}
\textbf{Calling the synchronous Textract API on multi-page PDFs.} It fails or truncates---and teams then split PDFs into single pages and call sync APIs per page, paying per-page latency and building a distributed fan-out by accident. Use the async API with SNS completion; that is what it exists for.
\end{pitfall}

\section{Tuesday: Comprehend and Translate}
\subsection{NLP as Enrichment}
Comprehend turns text into structured signal: entities (people, places, organizations), sentiment, key phrases, language, and---crucially for this book's governance thread---\emph{PII detection} with entity types and offsets \cite{awsComprehendDocs}. Custom classifiers and entity recognizers cover domain-specific needs when the general model is not enough. Translate handles language conversion at character-billed volume. Both integrate the same way: text in, JSON out, per-unit pricing---and both are happiest called in batches from a Lambda or Glue job rather than per row from a query engine.
\index{Comprehend}\index{Translate}

\begin{tipbox}
Comprehend's PII detection and Macie (Chapter~23) overlap but serve different moments: Comprehend is \emph{in-pipeline} (redact or route before persistence); Macie is \emph{at-rest discovery} (find what already landed). Mature designs use both.
\end{tipbox}

\section{Wednesday: Integrating AI APIs into Data Pipelines}
\subsection{The Integration Pattern}
The canonical shape: \textbf{S3 event $\rightarrow$ Lambda (or Glue job) $\rightarrow$ AI API $\rightarrow$ structured JSON to S3 $\rightarrow$ catalog + index}. Lambda suits per-object, sub-15-minute work; Glue suits bulk backfills and joins against existing tables. Three disciplines matter at scale:

\textbf{Idempotency and cost.} API calls are billed per unit; a retried pipeline that re-analyzes every object pays twice. Persist results keyed by object etag/version and skip already-processed inputs---Chapter~10's bookmark discipline applied to APIs.

\textbf{Result persistence as a dataset.} Extraction results are data: write them as versioned JSON under a governed prefix, register them in the catalog, and treat their schema as a contract (Chapter~11). ``The JSON from Textract'' becomes a table analysts join against everything else.

\textbf{Indexing for discovery.} Extracted metadata earns its keep when searchable: an OpenSearch domain indexed from the results stream (via Lambda or Firehose, Chapter~14) turns ten million processed documents into a queryable corpus \cite{awsOpenSearchDocs}.
\index{OpenSearch}

\subsection{PII in OCR Output}
OCR output is faithful: if the scanned form contained a national ID number, the JSON contains it too. Persisting raw extraction JSON is persisting PII. The pipeline must decide per field: keep (governed zone, masked grants), redact (Comprehend PII + substitution), or route to quarantine (Chapter~23's workflow). The default for raw outputs is the restricted zone with a TTL; curated, reviewed extracts are what analysts see.

\section{Thursday Hands-on Project: S3-Triggered Textract Extraction}
\index{hands-on project!Textract extraction}
Write a Lambda function triggered by an S3 upload that sends a scanned PDF to Textract, extracts key--value pairs, and saves them as JSON back to S3.

\subsection*{Lab Objectives}
\begin{itemize}
    \item Wire an S3 event notification to Lambda on a \texttt{raw/documents/} prefix.
    \item Call Textract (async for multi-page), parse the block structure into key--value JSON.
    \item Persist results idempotently under \texttt{extracted/documents/}.
    \item Handle a failure path visibly.
\end{itemize}

\subsection*{Procedure}
For the lab's single-page forms the synchronous API suffices; the code notes where the async boundary sits.

\begin{lstlisting}[language=Python,caption={S3-triggered Lambda extracting key--value pairs with Textract.},label={lst:ch20-textract}]
import json
import os
from urllib.parse import unquote_plus

import boto3

s3 = boto3.client("s3")
textract = boto3.client("textract")
OUT_PREFIX = os.environ.get("OUT_PREFIX", "extracted/documents/")

def lambda_handler(event, context):
    for rec in event["Records"]:
        bucket = rec["s3"]["bucket"]["name"]
        key = unquote_plus(rec["s3"]["object"]["key"])

        # Idempotency: skip if already extracted for this exact object version.
        out_key = f"{OUT_PREFIX}{key.rsplit('/', 1)[-1]}.json"
        head = s3.head_object(Bucket=bucket, Key=key)
        etag = head["ETag"].strip('"')
        try:
            existing = s3.head_object(Bucket=bucket, Key=out_key)
            if existing.get("Metadata", {}).get("source-etag") == etag:
                continue
        except s3.exceptions.ClientError:
            pass

        # Sync API: single-page documents only. Multi-page -> StartDocumentAnalysis
        # with SNS completion and a second Lambda handling the result.
        resp = textract.analyze_document(
            Document={"S3Object": {"Bucket": bucket, "Name": key}},
            FeatureTypes=["FORMS"],
        )

        # Reconstruct key-value pairs from the block graph.
        blocks = {b["Id"]: b for b in resp["Blocks"]}
        kv = {}
        for b in resp["Blocks"]:
            if b["BlockType"] == "KEY_VALUE_SET" and "KEY" in b["EntityTypes"]:
                key_text = _text_of(b, blocks)
                val_text = ""
                for rel in b.get("Relationships", []):
                    if rel["Type"] == "VALUE":
                        for vid in rel["Ids"]:
                            vb = blocks[vid]
                            val_text = _text_of(vb, blocks)
                if key_text:
                    kv[key_text] = val_text

        s3.put_object(
            Bucket=bucket, Key=out_key,
            Body=json.dumps({"source": key, "fields": kv}, indent=2).encode(),
            ContentType="application/json",
            Metadata={"source-etag": etag},
        )

def _text_of(block, blocks):
    parts = []
    for rel in block.get("Relationships", []):
        if rel["Type"] == "CHILD":
            for cid in rel["Ids"]:
                child = blocks[cid]
                if child["BlockType"] == "WORD":
                    parts.append(child["Text"])
    return " ".join(parts)
\end{lstlisting}

\subsection*{Verification}
\begin{itemize}
    \item Uploading a sample form produces a JSON file with the expected key--value pairs.
    \item Re-uploading the identical object does not re-call Textract (idempotency via the etag metadata check).
    \item A corrupted file produces a visible failure (Lambda error metric, retry, then a documented dead-letter outcome).
    \item The output prefix is registered in the catalog and queryable in Athena.
\end{itemize}

\section{Friday Use Case: Terabyte-Scale Image Understanding with Search}
\index{use case!image processing pipeline}
A photo-sharing platform receives terabytes of user images daily. Product wants every image searchable by its content (``receipts,'' ``dogs on beaches,'' ``documents with signatures''), and trust-and-safety wants moderation at upload. Design the pipeline.

\subsection{Architecture}
\begin{figure}[H]
    \centering
    \resizebox{\textwidth}{!}{%
    \begin{tikzpicture}[
        box/.style={rectangle, rounded corners, draw=primary, thick, fill=primary!6, align=center, minimum width=2.9cm, minimum height=1.0cm, font=\small\bfseries},
        storebox/.style={cylinder, shape aspect=0.25, draw=primary, thick, fill=blue!6, shape border rotate=90, align=center, text width=2.5cm, minimum height=1.25cm, font=\footnotesize\bfseries},
        guard/.style={rectangle, rounded corners, draw=red!60!black, thick, fill=red!5, align=center, minimum width=2.6cm, minimum height=0.9cm, font=\footnotesize\bfseries},
        arr/.style={-{Stealth[length=2.5mm]}, draw=primary, thick}
    ]
        \node[storebox] (raw) at (0,0) {S3 uploads\\(raw images)};
        \node[box] (fn) at (3.2,0) {Lambda\\orchestrator};
        \node[box] (rek) at (6.4,1.1) {Rekognition\\labels, moderation};
        \node[box] (tex) at (6.4,-1.1) {Textract\\embedded text};
        \node[storebox] (meta) at (9.8,0) {S3 metadata\\(cataloged JSON)};
        \node[box] (os) at (13.0,0) {OpenSearch\\searchable index};
        \node[guard] (mod) at (9.8,2.2) {Moderation queue\\(confidence-gated)};
        \draw[arr] (raw) -- (fn);
        \draw[arr] (fn) -- (rek);
        \draw[arr] (fn) -- (tex);
        \draw[arr] (rek) -- (meta);
        \draw[arr] (tex) -- (meta);
        \draw[arr] (meta) -- (os);
        \draw[arr, dashed] (rek) -- (mod);
    \end{tikzpicture}%
    }
    \caption{Image understanding pipeline: parallel enrichment by Rekognition and Textract, cataloged JSON metadata, OpenSearch indexing, and a confidence-gated moderation path.}
    \label{fig:ch20-image-pipeline}
\end{figure}

\subsection{Design Decisions}
\textbf{Fan-out per image.} One S3 event triggers a Lambda that invokes Rekognition (labels + moderation) and, only for images with detected text density, Textract---conditional invocation is a cost decision at terabyte scale. Results converge into one metadata document per image.

\textbf{Confidence thresholds as policy.} Labels below 80\% confidence are stored but not searchable; moderation scores above the action threshold route to a human-review queue with the image reference---never auto-deletion on a model's say-so. Thresholds are configuration, versioned, and tuned from review-outcome statistics.

\textbf{The metadata is the product.} Cataloged JSON (labels, OCR text, moderation verdict, EXIF) in the governed lake feeds analytics; an OpenSearch index fed from the same documents serves interactive search. Two consumers, one extraction: the write-once, read-many discipline of the whole book applied to AI output.

\textbf{Cost shape.} Per-image Rekognition calls dominate; Textract runs on a filtered subset; OpenSearch sizing follows index volume and query rate. The monthly model states per-1{,}000-image costs and the levers (label-only tiers, text-density gating, index refresh interval).

\begin{pitfall}
\textbf{Persisting OCR text without a PII posture.} User images contain documents, ID cards, and screenshotted messages. The extracted text is PII-bearing data the moment it lands: govern the raw outputs, redact in the searchable index, and let Chapter~23's Macie scanning verify the posture continuously.
\end{pitfall}

\section{Interview Q\&A}
\subsection{Concepts and Trade-offs}
\begin{enumerate}
    \item \textbf{Q: Textract versus Comprehend: which extracts key--value pairs from scanned forms?}\\
    A: Textract---\texttt{AnalyzeDocument} with \texttt{FORMS} returns the form's key--value structure. Comprehend operates on already-digital text (entities, sentiment, PII), often as Textract's downstream.
    \item \textbf{Q: Why trigger AI APIs from S3 events via Lambda?}\\
    A: Arrival-driven processing with no polling and per-object scale-out; the event carries the object identity, and Lambda's concurrency model matches spiky upload patterns.
    \item \textbf{Q: When must you switch from synchronous to asynchronous Textract jobs?}\\
    A: For multi-page PDFs/TIFFs or large documents. The async API runs a job and signals completion via SNS; the sync API is single-page only.
    \item \textbf{Q: How do you make extracted document metadata searchable?}\\
    A: Persist normalized JSON results, then index them into OpenSearch (via Lambda or Firehose) while also registering the results in the catalog for SQL analytics.
    \item \textbf{Q: What PII risks exist when persisting raw OCR output?}\\
    A: OCR is faithful---IDs, account numbers, and addresses in the document appear in the JSON. Raw outputs need a restricted zone, field-level redaction for curated zones, and at-rest discovery scanning.
    \item \textbf{Q: How do you keep retried pipelines from re-billing API calls?}\\
    A: Idempotency keyed on object identity: persist results with the source etag/version and skip already-processed inputs.
\end{enumerate}

\section{Further Reading}
The Textract and Rekognition developer guides cover block structure, async jobs, and confidence semantics \cite{awsTextractDocs,awsRekognitionDocs}; the Comprehend guide covers NLP and PII detection \cite{awsComprehendDocs}; the OpenSearch guide covers indexing patterns \cite{awsOpenSearchDocs}. The event and Lambda foundations are in \cite{awsEventBridgeDocs,awsLambdaDocs}.

\section*{Key Takeaways}
\begin{itemize}
    \item Managed AI APIs are per-record enrichment transforms with a network boundary; pipeline discipline is the engineering content.
    \item The sync/async boundary (single-page versus multi-page) determines pipeline shape; async means event-driven completion.
    \item Persist extraction results as governed, cataloged datasets with idempotency keyed on object identity.
    \item Confidence thresholds are versioned policy, and moderation decisions stay human-gated.
    \item Raw OCR output is PII-bearing data: restrict it, redact for curated use, and scan continuously.
    \item One extraction, many consumers: catalog for analytics, OpenSearch for interactive discovery.
\end{itemize}

\section{Exercises}
\begin{enumerate}
    \item \textbf{(Conceptual)} Explain why an asynchronous document API forces an event-driven pipeline design.
    \item \textbf{(Conceptual)} Why store low-confidence labels but exclude them from search, rather than discarding them?
    \item \textbf{(Hands-on)} Run the Thursday lab; then convert a multi-page PDF path to the async API with SNS completion.
    \item \textbf{(Hands-on)} Add Comprehend PII detection to the extraction Lambda and redact detected spans in the curated output.
    \item \textbf{(Hands-on)} Index the extracted JSON into a small OpenSearch domain and query it.
    \item \textbf{(Design)} Design the moderation-review workflow: queueing, reviewer tooling inputs, decision persistence, and threshold tuning feedback.
    \item \textbf{(Design)} Build the monthly cost model for the Friday pipeline at 50 million images/day; identify the top two cost levers.
    \item \textbf{(Architecture)} Extend the architecture for video uploads: which Rekognition APIs apply, and what changes about latency and cost?
\end{enumerate}

\section{Milestone 5 Capstone: End-to-End MLOps Pipeline}\label{sec:ch20-capstone}
\index{capstone project}\index{Milestone 5 capstone}\index{MLOps!capstone}

\begin{capstonebox}
\textbf{Mission.} Deliver the full MLOps lifecycle orchestrated by MWAA: Athena extracts features into SageMaker Feature Store, a custom training job produces a model, the Model Registry gates it, and Model Monitor watches the serving endpoint for training/serving skew---with MWAA owning the schedule and the alerts.

\textbf{Estimated time.} 8--10 hours in a sandbox AWS account.

\textbf{What you\textquotesingle{}ll practice.}
\begin{itemize}
    \item Feature engineering from the lake (Athena) into Feature Store (Weeks 18--19).
    \item Custom SageMaker training with versioned inputs.
    \item Registry gating, monitored deployment, and drift feedback.
    \item Airflow orchestration of the whole lifecycle (Week~12).
\end{itemize}
\end{capstonebox}

\subsection*{Scenario}
The retailer's personalization team ships a propensity model monthly. The platform team must turn that into infrastructure: one DAG run produces features, a trained model, an evaluated registry entry, and a monitored deployment---with serving skew watched continuously.

\subsection*{Requirements}
\textbf{Functional requirements.}
\begin{itemize}
    \item An Athena feature extraction writing to a feature group (offline for training, online for serving).
    \item A SageMaker Pipeline (preprocess $\rightarrow$ train $\rightarrow$ evaluate $\rightarrow$ register) with a metric gate.
    \item MWAA orchestration from feature extraction through pipeline execution.
    \item A monitored endpoint with a skew/violation alarm to SNS and an MWAA failure alert.
\end{itemize}

\textbf{Non-functional requirements.}
\begin{itemize}
    \item End-to-end runs from a clean environment per the README.
    \item At least three automated tests: feature validation, a model-quality gate, and a drift-injection test demonstrating a Monitor violation.
    \item A monthly cost estimate (training instance-hours, endpoint versus Batch Transform, Feature Store reads/writes) with two documented optimizations.
    \item Security review: training access scoped via Lake Formation, artifacts in a KMS-encrypted bucket, no credentials in notebooks.
\end{itemize}

\subsection*{Reference Architecture}
The pipeline follows \cref{fig:ch19-mlops-loop}: Athena extraction $\rightarrow$ Feature Store $\rightarrow$ Pipeline $\rightarrow$ Registry $\rightarrow$ monitored endpoint, with MWAA orchestrating and drift violations feeding back through EventBridge to retraining with human approval.

\subsection*{Implementation Steps}
\begin{enumerate}
    \item Create the feature group and the Athena-to-Feature-Store ingestion step.
    \item Define the pipeline with the quality gate (F1 $\geq$ 0.80) and registry integration.
    \item Author the MWAA DAG: extraction $\rightarrow$ pipeline start $\rightarrow$ status watch $\rightarrow$ notify.
    \item Deploy the approved model with data capture; build the baseline and monitoring schedule.
    \item Implement the three tests, including drift injection producing a real violation.
    \item Write the cost model and the security review.
\end{enumerate}

\subsection*{Deliverables}
\begin{itemize}
    \item The DAG, pipeline definition, and training code in version control.
    \item Test evidence, including the injected drift violation and the gate rejection of a weak model.
    \item Alarm configurations and one fired alert of each kind.
    \item The cost model with two optimizations and the security review.
\end{itemize}

\subsection*{Acceptance Criteria}
\begin{itemize}
    \item One DAG run carries a dataset from extraction to a gated, registered, deployable model.
    \item A weak model is rejected at the gate; a drift event reaches SNS with context.
    \item Training data access is demonstrably catalog-governed; artifacts are KMS-encrypted.
    \item Endpoint shape is justified against the traffic estimate (serverless or batch by default).
\end{itemize}

\subsection*{Stretch Goals}
\begin{itemize}
    \item Add a champion/challenger shadow deployment with live comparison metrics.
    \item Wire the Week~20 enrichment outputs as additional features and measure lift.
    \item Add lineage emission (OpenLineage) from the DAG and pipeline.
    \item Automate monthly baseline refresh tied to the retraining cadence.
\end{itemize}
